\section{Formal Languages} 

\subsection{Alphabets, Strings, and Languages}

	\begin{definition}[Alphabet]
		\label{def:alphabet}
    An \textbf{alphabet} $\Sigma$ is a set. 
  \end{definition}

  \begin{definition}[Kleene Star]
    Given a set $\Sigma$, we define
    \begin{equation}
      \Sigma^\ast \coloneqq \bigcup_{n=0}^\infty \Sigma^n
    \end{equation}
    Each element of $\Sigma^\ast$ is called a \textbf{word} or a \textbf{string}, and a subset of $\Sigma^\ast$ is called an \textbf{expression}. The \textbf{empty word}, denoted $\epsilon$, is the unique string of length $0$. 
  \end{definition}

  Now let's define some operations. 

  \begin{definition}[Concatenation]
    Given two words $u, v \in \Sigma^\ast$, the \textbf{concatenation} $uv = u \cdot v$ is the word formed by appending the sequence of symbols in $v$ to the sequence of symbols in $u$. 
  \end{definition}

  Note that $(\Sigma^\ast, \cdot, \epsilon)$ is a monoid, where $\epsilon$ represents the empty word. Some strings (programs) are legal while others are not. Therefore, we can define the set of all valid programs as simply as a subset of $\Sigma^\ast$. 

  \begin{definition}[Formal Language]
    A \textbf{formal language} over alphabet $\Sigma$ is a subset $L \subset \Sigma^\ast$. 
    \begin{enumerate}
      \item A word $w \in \Sigma^\ast$ is \textbf{well-formed} if $w \in L$. 
      \item An expression $E \subset \Sigma^\ast$ is \textbf{well-formed} if $E \subset L$. 
    \end{enumerate}
  \end{definition}

	We can do operations on languages. 

  \begin{definition}[Product of Formal Languages]
    Given two formal languages $L_1, L_2 \subset \Sigma^\ast$, their product is defined 
    \begin{equation}
      L_1 L_2 \coloneqq \{ uv \mid u \in L_1, v \in L_2 \}
    \end{equation}
  \end{definition}

\subsection{Decision Problems as Languages}

	Decision problems as languages

\subsection{Encodings}

	Encodings of mathematical objects as strings

\subsection{Formal Grammars}

  Generally, a subset of $\Sigma^\ast$ doesn't give us much structure, so we would like to define some way to pick the subset $L$ out. We can do this with \textit{grammars}. 

  \begin{definition}[Formal Grammar]
    A \textbf{formal grammar} is a 4-tuple $G = (V, \Sigma, R, S)$ consisting of the following. 
    \begin{enumerate}
      \item \textit{Non-terminal Symbols}. $V$ is a finite set consisting of \textbf{non-terminal symbols}, also known as \textbf{variables}. 
      \item \textit{Terminal Symbols}. $\Sigma$ is a finite set---disjoint from $V$---consisting of \textbf{terminals}. 
      \item \textit{Production Rule}. $R$ is a relation, i.e. is a finite subset of 
        \begin{equation}
          (V \cup \Sigma)^\ast V (V \cup \Sigma)^\ast \times (V \cup \Sigma)^\ast
        \end{equation}
        where the LHS represents the product of the languages. We usually write the relation as $\alpha \to \beta$.\footnote{In math, this is usually written $\alpha R \beta$, or $\alpha \sim \beta$ if it is an equivalence relation.}

      \item \textit{Start Symbol}. $S \in V$ is the initial variable from which derivation begins. 
    \end{enumerate}
  \end{definition}

  We can think of the production rule as a set of transformations---formally called \textit{derivations}---that you can apply to a word. 

  \begin{definition}[Derivation]
    Let $G = (V, \Sigma, R, S)$ be a formal grammar. We define a binary relation $\derives$ on the set of all strings $(V \cup \Sigma)^\ast$ as follows: 
    \begin{enumerate}
      \item A string $u$ \textbf{directly derives} $v$, written $u \derives v$, if there exists strings $\phi, \psi \in (V \cup \Sigma)^\ast$ and a production rule $(\alpha \to \beta) \in R$ such that 
        \begin{equation}
          u = \phi \alpha \psi, \qquad v = \phi \beta \psi
        \end{equation}
        In other words, we say a word $w$ is derived from $v$, written $v \derives w$, if $w$ can be obtained by replacing a part of $v$ according to a rule in $R$. 

      \item We say $u$ \textbf{derives} $v$\footnote{We also say that $v$ is in the transitive closure of $u$.}, written $u \tderives v$, if $u$ can be directly derived into $v$ in $0$ or more steps. The relation $\tderives$, called the \textbf{reflexive transitive closure} of $\derives$, is defined as follows. $u \tderives v$ if 
        \begin{enumerate}
          \item $u = v$, or 
          \item There exists a finite sequence of strings $u = w_0, w_1, \ldots, w_n = v$ such that $w_i \derives w_{i+1}$ for all $0 \leq i < n$. 
        \end{enumerate}

    \end{enumerate}
  \end{definition}

  \begin{definition}[Language Derived From Grammar]
    The \textbf{language derived from $G$}, denoted $L(G)$, is the set of all terminal strings that can be derived from $S$ using the rules in $R$. It is defined: 
    \begin{equation}
      L(G) \coloneqq \{ w \in \Sigma^\ast \mid S \derives^\ast w \}
    \end{equation}
  \end{definition}

  \begin{example}
    Let $V = \{E\}$ be the non-terminals, and $\{+, -, \texttt{num}\}$ be our terminals, which come from our lexer. Then, our production rules can look something like. 
    \begin{enumerate}
      \item $E \to E + E$. 
      \item $E \to E - E$. 
      \item $E \to \texttt{num}$. 
    \end{enumerate}
    So basically, if we have some word, then we can replace any $E$ in the word by any of the three choices on the right hand side. For example, we can do the following sequence of derivations. 
    \begin{align}
      E & \derives E + E \\ 
        & \derives E - E + E \\
        & \derives E - \texttt{num} + E\\
        & \derives E + E - \texttt{num} + E 
    \end{align}
    You can keep doing this until there is no more production rules you can apply. For context free grammars, you basically do this until there are only terminals left, e.g. $\texttt{num} + \texttt{num} - \texttt{num} + \texttt{num}$. 
  \end{example}

	You can track the derivations used to go from the base expression to the sequence of terminals. This is perfectly represented by a \hyperref[dsa-def:tree]{tree}. 

  \begin{definition}[Parse Tree]
    Given a formal grammar $G = (V, \Sigma, R, S)$, a \textbf{parse tree} is an ordered tree where 
		\begin{enumerate}
		  \item the root is the start symbol $S$. 
			\item each internal node is a nonterminal symbol. 
			\item each leaf is either a terminal symbol or $\epsilon$ 
			\item if an internal node is labeled by a nonterminal $A$, and its children from left to right are labeled $X_1, \ldots, X_k$, then the grammar must contain the production 
				\begin{equation}
				  A \derives X_1 X_2 \ldots X_k
				\end{equation}
		\end{enumerate}

		The \textbf{yield} of a parse tree is the string obtained by reading the leaves from left to right. 
  \end{definition}

  \begin{definition}[Transitive Closure]
  \end{definition}

  \begin{definition}[Language Generated by Grammar]
    Given grammar $G = (V, \Sigma, R, S)$, the language $L(G)$ is the set of all terminal strings that can be derived from the start symbol $S$. 
    \begin{equation}
      L(G) \coloneqq \{ w \in \Sigma^\ast \mid S \derives^\ast w \}
    \end{equation}
  \end{definition}

  Therefore, you can basically think of the ``complexity'' or size of a language $L$ as being determined by the ``complexity'' of the generating grammar $G$. Usually, the alphabet is kept fixed, and the main contributor to the complexity are the production rules $R$. Depending on how much we restrict $R$, we get different levels of languages in the \textit{Chomsky Hierarchy}. 

  \begin{definition}[Unrestricted]
    An \textbf{unrestricted language} $L$ is a language that can be generated by some grammar $G$. 
  \end{definition}

	As the name suggests, there will be languages which are restricted. The hierarchy of these restricted languages is called the \textbf{Chomsky hierarchy}. 

